%==============================================================================
%  PAPER 2 -- Appendix
%
%  Supporting figures use non-floating column-width blocks so captions remain
%  in reading order and final-page balancing cannot reorder pending floats.
%==============================================================================

\section{Source Code and Reproduction}
\label{app:t2-code}

The complete implementation --- the controller, the genetic tuning experiments,
the benchmark study, the MATLAB port and every script that produced a figure or
a number in this paper --- is available at~\cite{repo2}:

\begin{center}
\small\url{https://github.com/bpdyl/ai-research-topic-modeling}
\end{center}

\noindent Task 2 lives under \texttt{assistive-care-flc/} and reproduces from a
clean checkout. Three drivers regenerate the experiments in order:
\texttt{run\_part1.py} (the controller and its structural audit),
\texttt{run\_part2.py} (genetic tuning with its budget-matched control) and
\texttt{run\_part3.py} (the 180 benchmark runs). A fourth,
\texttt{make\_report\_figures.py}, derives the remaining figures and the two
audits of Section~\ref{sec:t2-cost} from the stored results without re-running
any experiment.

The MATLAB system is generated from the controller definition used by
the reference engine; parity is checked on selected examples rather than assumed. It requires no
Fuzzy Logic Toolbox and runs on base MATLAB. All random seeds are fixed, and an
automated test suite pins the behavioural claims made in this paper, including
the safety property that the controller never heats a hot room.

\section{Supporting Figures}
\label{app:t2}

The figures below support claims made in the body of the paper. Each one
substantiates a statement the body makes in a sentence, and is placed here so
that the argument reads continuously without a full-page figure interrupting
it. Every figure is regenerated by the scripts listed in
Appendix~\ref{app:t2-code} from stored results, with the benchmark panels regenerated
for the official-shift experiment.


\begin{appendixfigure}
\centering
\includegraphics[width=\columnwidth]{fig05_inference_stages_dimmer.png}
\caption{Implication, aggregation and defuzzification for the \texttt{dimmer}
output on the worked scenario --- the counterpart to Fig.~\ref{fig:t2-stages},
which shows the same three stages for \texttt{hvac}. Two lighting rules fire:
\texttt{Dark} $\wedge$ \texttt{Active} $\rightarrow$ \texttt{High} at
$\alpha = 0.25$, and \texttt{Dim} $\wedge$ \texttt{Active} $\rightarrow$
\texttt{Medium} at $\alpha = 0.20$. Each consequent set is clipped at its own
$\alpha$ (the flat tops), the two are unioned, and the centroid of the result
is 72.60\%. Note that the higher-activation rule is the one demanding
\emph{more} light, which is why the crisp result sits well above the midpoint
of the two sets rather than between them.}
\label{fig:t2-stages-dimmer}
\end{appendixfigure}

\begin{appendixfigure}
\centering
\includegraphics[width=\columnwidth]{fig07_surface_hvac_temp_preference.png}
\caption{Control surface: \texttt{hvac} against \texttt{room\_temp} and
\texttt{preference}, at light activity. This is the slice that shows the
clamped modifier of \eqref{eq:modifier} doing its work. Preference shifts the
surface bodily up and down --- a resident who prefers warmth is given a warmer
room at every temperature --- but it never inverts the slope in temperature. A
resident who sets the dial to $+5$ therefore cannot cause a hot room to be
heated, which is the safety property the inner clip guarantees by
construction.}
\label{fig:t2-surface-pref}
\end{appendixfigure}

\begin{appendixfigure}
\centering
\includegraphics[width=\columnwidth]{fig09_response_hvac_by_activity.png}
\caption{HVAC command against room temperature at three activity levels, at
neutral preference. Three things are visible. The curves are ordered as the
design intends: a resting occupant is given a warmer room than an active one at
every temperature, which is the metabolic-heat argument of
Section~\ref{sec:t2-intro} realised in the rule base. The response saturates
symmetrically at roughly $\pm 79.9\%$, the centroid limit quantified in
Section~\ref{sec:t2-behaviour}. And the small non-monotonicity noted there ---
about 0.6\% of the output span --- is visible as a shallow reversal near the
points where the demand scale saturates and two adjacent temperature sets come
to share a consequent.}
\label{fig:t2-response}
\end{appendixfigure}


\begin{appendixfigure}
\centering
\includegraphics[width=\columnwidth]{fig13_pred_vs_target.png}
\caption{Controller output against the reference policy on the 120-example
held-out split, before and after tuning, for both actuators. The tuned points
lie visibly closer to the diagonal. The residual spread is expected rather than
a failure: the reference policy is a saturating linear ramp and the controller
is a union of clipped fuzzy sets, so the two functional forms cannot coincide
exactly. That mismatch is deliberate --- had the target been generated by a
Mamdani system with the same rule base, the error could in principle have been
driven to zero, which would have demonstrated only that the encoding can
represent its own output.}
\label{fig:t2-predtarget}
\end{appendixfigure}

\begin{appendixfigure}
\centering
\includegraphics[width=\columnwidth]{fig23_encoding_cost.png}
\caption{What the 8-bit-per-gene binary encoding costs before the search
begins. \textbf{Left:} the round-trip displacement of each of the 81
membership-function parameters after encoding the hand-designed controller and
decoding it again --- at most 0.392 units, which is half a quantisation step on
the widest universe. \textbf{Right:} the training error that displacement
causes, $+0.000342$ in normalised RMSE, set against the improvement the search
subsequently wins. The encoding costs roughly 0.7\% of what the search gains,
which is the evidence for the claim in Section~\ref{sec:t2-ga} that
quantisation is not what limits the result.}
\label{fig:t2-encoding}
\end{appendixfigure}


\begin{appendixfigure}
\centering
\includegraphics[width=\columnwidth]{fig21_benchmark_landscapes.png}
\caption{The two CEC'2005 instances actually optimised, at $D=2$, plotted as
$\log_{10}(f - f^{*})$ so that structure is visible on both scales at once.
F6's long, narrow, curved valley and F9's regular lattice of local minima are
the complementary difficulties described in Section~\ref{sec:t2-cec}: one is
hard because the descent direction keeps turning, the other because there are
$10^{D}$ basins to be trapped in. The marked optima confirm the shift vectors
are doing their job --- neither optimum sits at the centre of its domain, so an
algorithm that initialises symmetrically about the origin gains no free
advantage. The shift vectors are taken from the official CEC2005 distribution.}
\label{fig:t2-landscapes}
\end{appendixfigure}

\balance
\section{MATLAB benchmark implementation}
The following function implements both selected benchmark equations using the
vendored official shift files. Its Python counterpart was checked against the
same published data and known optima. This additional MATLAB listing was not
executed in the report-revision environment.
\lstinputlisting[language=Matlab,basicstyle=\ttfamily\scriptsize,breaklines=true,columns=fullflexible]{code/cec2005_f6_f9.m}
